\chapter*{结\quad 论}
\addcontentsline{toc}{chapter}{结\quad 论}

本文设计并实现了一套基于树莓派的人脸识别考勤系统。系统面向实验室门口、小型办公场所等低成本考勤场景，以树莓派作为边缘终端，完成了从成员名单导入、用户注册、人脸采集、管理员审核、终端识别到自动签到记录的完整业务闭环。

本文的主要工作如下。

第一，完成了完整的用户注册与人脸资料提交流程。系统通过本地CSV名单限制允许注册的用户范围，要求用户在首次注册时完成姓名、编号和初始密码校验，并强制设置正式密码。人脸采集环节采用浏览器摄像头现场抓拍，结合后端动作挑战会话，保证最终提交的注册照来自通过挑战的可信抓拍结果。

第二，完成了注册端人脸质量校验与审核机制。系统设计了可扩展的校验流水线，包含单人检测、人脸尺寸、模糊度、亮度、关键点完整性、姿态和翻拍风险检测。通过校验的照片进入待审核状态，只有管理员审核通过后才会进入终端识别人脸库。驳回记录会保存原因和备注，供用户重新上传时参考。

第三，完成了终端实时识别与自动签到。系统采用OpenCV YuNet进行人脸检测，采用SFace进行特征提取与余弦相似度匹配，结合单人入镜、人脸面积、活体检测、身份匹配、考勤时间窗和重复签到冷却等门禁，避免错误签到和重复签到。签到成功后，系统写入SQLite考勤记录并保存现场快照。

第四，完成了被动活体检测与重放风险分析。系统通过OpenVINO加载anti-spoof-mn3模型，对终端人脸区域进行真人概率判断，并使用2秒滑动窗口融合多帧结果。同时，系统利用反光、边缘、频域和矩形边框等轻量信号估计屏幕重放风险，在不增加用户额外动作负担的情况下增强终端识别安全性。

第五，完成了远程管理接口设计。树莓派端提供\texttt{/api/admin/*}管理员API，支持设备信息、健康状态、性能指标、成员管理、用户管理、人脸审核、考勤查询和快照查看。Windows端管理系统可以通过Tailscale网络访问这些接口，实现远程审核和运维。

系统目前仍存在一些不足。首先，系统主要面向小规模人脸库，在人数进一步扩大时，线性相似度匹配的效率可能下降，后续可以引入向量索引加速检索。其次，当前活体检测主要依赖RGB图像和轻量风险信号，对于高质量屏幕重放攻击仍有提升空间，后续可以探索更多时序特征和生理信号。第三，系统目前以单台树莓派为核心，尚未实现多终端统一管理和跨设备数据同步。第四，用户门户和管理员页面仍有继续优化空间，可以进一步提升交互体验和统计分析能力。

未来工作可以从以下方向展开：接入可配置的考勤规则引擎，支持迟到、早退、请假和补签等业务；扩展多设备协同能力，实现多个树莓派终端的统一管理；增加考勤数据导出、统计报表和可视化分析；优化活体检测策略，提升对复杂翻拍和预录视频攻击的识别能力；在大规模人脸库场景下引入高效向量检索方法，进一步提升系统性能。
